% Chapter 4, Figure 16 (scan: figures/ch4/fig16.png): a Malewicki chart, y_max of vertical flight against the
% liftoff mass m_o for B4-powered models, k = .00005 ... .0016 kg/m lettered on the curves (Bengen's maxima).
% Curves: fig16.py, COMPUTED by the Fehskens-Malewicki solution, eqs. (20), (21), with the coast of eq. (67)
% (trajectory.fehskens_malewicki: mean mass, F = I_t/t_b; B4 from the Figure 4 lettering, propellant 8.33 g from
% Table 1, g = 9.8), from the engine alone (m_o = .021 kg, Fig 11's curve (a)) to .10 kg; the line of Bengen's
% maxima is the locus of the curves' maxima as k varies, from 700 m (k = .000032) down to the .0016 curve, where
% the printed line ends. Overlay on the scan: the six curves 95% within 1-1.4 px (0.24 mm) of the printed ones;
% the Bengen line within 3.6 px (0.61 mm), the computed maxima of the high-k curves lying left of the printed
% line's lower end (k = .0016: .0378 kg against about .040 drawn). The book's interval method, eqs. (83)-(87),
% is METHOD = "interval" in fig16.py (curves 1-4 m away, lower end of the Bengen line at .0407 kg); that switch
% changes the CSVs only, and the k labels and the callout below, placed for these curves, would have to be moved
% (the Bengen line would then cross ".0016" and run along ".0008" and ".0004"). For that variant (checked in a
% render): .0001 to x = .0293, .0002 .0330, .0004 .0365, .0008 .0396, .0016 .0412 (y unchanged; .00005 stays),
% the callout point to (.0240, 640) and its elbow to (.0293, 668). The k labels are \footnotesize, as are the
% curve labels and legends of the other Chapter 4 plots; the callout and the engine label stay \small.
% A design chart read for values: hairline grid at the printed spacing (.01 kg, 100 m). The six curves are one
% family in one colour (s1) labelled on the curves as printed: just above each maximum, right of the Bengen line
% (.00005 beside its falling branch); the Bengen line s2 dashed (printed long-dash-short-dash); the engine's own
% mass a reference line (guide; printed short-dashed), its label set along it (the 1973 label stands further
% left with a short leader stroke across the line). Labels in white knock-outs where the 1973 art breaks the grid.
\documentclass[11pt]{standalone}
\usepackage{tamrfig}
\begin{document}
\begin{tikzpicture}
  \begin{axis}[tamr grid, grid=both,
      width=3.8in, height=5.2in,
      xmin=0, xmax=0.10, xtick={0,0.02,...,0.10}, xticklabels={0, 0.02, 0.04, 0.06, 0.08, 0.10},
      minor xtick={0.01,0.03,...,0.09},
      ymin=0, ymax=700, ytick distance=100,
      xlabel={$m_o$ (kg)}, ylabel={$y_{\max}$ (m)}]
    % the mass of the loaded engine alone: no model can be lighter
    \draw[guide] (axis cs:0.021,0) -- (axis cs:0.021,700);
    \node[rotate=90, anchor=south, fill=white, inner sep=1.5pt, yshift=1pt] at (axis cs:0.021,300) {$m_o$ of engine alone};
    % y_max(m_o) for each k (kg/m)
    \pgfplotsinvokeforeach{k00005,k0001,k0002,k0004,k0008,k0016}{
      \addplot[s1] table[x=m0, y=#1, col sep=comma] {figures/v2/ch4/fig16.csv};
    }
    % the locus of the maxima
    \addplot[series2] table[x=m0, y=ymax, col sep=comma] {figures/v2/ch4/fig16-bengen.csv};
    \begin{scope}[every node/.style={fill=white, inner sep=1.2pt, font=\footnotesize}]
      \node[anchor=west]       at (axis cs:0.0400,468) {.00005};
      \node[anchor=south west] at (axis cs:0.0288,358) {.0001};
      \node[anchor=south west] at (axis cs:0.0322,237) {.0002};
      \node[anchor=south west] at (axis cs:0.0354,158) {.0004};
      \node[anchor=south west] at (axis cs:0.0379,106) {.0008};
      \node[anchor=south west] at (axis cs:0.0384,72) {.0016};
    \end{scope}
    \coordinate (bengen) at (axis cs:0.02375,640);
    \coordinate (elbow) at (axis cs:0.0290,668);
  \end{axis}
  % (\callout measures its points, so it is drawn after the axis)
  \tikzset{callout label/.append style={fill=white}}
  \callout{(bengen)}{(elbow)}{Line of Bengen's maxima}
\end{tikzpicture}
\end{document}
